\section{Design Rationale: Alternatives Considered and Rejected}
\label{app:rationale}

Section~\ref{ch:model} states the rules the model uses. This appendix records the alternatives
considered and rejected for each.

\subsection{Population Construction}

\textbf{Statistical fusion of \ac{NIDS} and \ac{FinScope}.} A hot-deck or model-based fusion would
recover finer joint structure between the two instruments than copying flags within an
income-quintile-by-province cell, at the cost of its own validation burden and a real risk of
fitting the idiosyncrasies of two surveys rather than the population behind them. The crude match
reproduces cell marginals, which is what the model needs.

\textbf{\ac{CPI} forwarding to 2022.} An earlier design forwarded \ac{NIDS} monetary values to
2022, which required a second reference year, assumptions about how relative incomes moved between
the two, and 2022-vintage validation targets. Working entirely in 2017 Rands removes all three, and
the 2017 anchor is independently required by the counterfactual injection design.

\subsection{Household Decision Rules}

\textbf{A propensity-to-consume or buffer-stock consumption rule (Submodel~2).} Both require a
calibrated propensity for which no South African estimate exists, and both discard the
committed-versus-discretionary expenditure split already observed in the survey data.

\textbf{Cost ranking as the lender-choice rule (Submodel~5).} \ac{BNPL} is nominally interest-free,
so a cost-minimising rule selects it trivially and attributes the choice to price. The evidence
attributes the choice to convenience and social approval \cite{ackert2025bnpl}, and the two
mechanisms imply different interventions: a price-driven choice would respond to disclosure of
effective cost, while a norm-driven one would not.

\textbf{Scheduled amortisation as the sole repayment rule (Submodel~6).} Keys and Wang
\cite{Keys2019} find that 29\% of credit-card accounts pay at or near the contractual minimum, and
that near-minimum payers shift their payments when the formula changes, which is the signature of
anchoring. A uniform amortisation rule would suppress the behaviour by which obligations persist
and compound.

\textbf{A debt-service-to-income threshold as the default condition (Submodel~7).} Any such
threshold is arbitrary and the \ac{NCA} prescribes none. Madeira's cash-flow insolvency test is
cited, consistent with the affordability rule the model already implements, and validated against
observed default rates in a middle-income economy \cite{madeira2018chile}.

\textbf{A continuous present-bias parameter (Submodel~8).} No South African distribution exists
against which to calibrate one, so it would add a free parameter with no discipline. The model
carries a single binary behavioural type instead, whose share is sourced and swept.

\subsection{Lender and Platform Rules}

\textbf{A flat \ac{DSTI} cap as the credit-granting gate (Submodel~9).} An earlier version of this
model used a ceiling of 0.65. It was unsourced, the \ac{NCA} prescribes no such ratio, and it was
badly over-permissive at the bottom of the distribution, since a 65\% rule would permit a household
earning R900 per month to service R585 per month. The Regulation 23A residual-income test produces
an income-varying ceiling, which across the population spans zero to 93.2\% of income.

\textbf{Making the \ac{BNPL} eligibility screen a real affordability test (Submodel~12).} The
asymmetry between the Regulation 23A test that binds registered lenders and the light screen
\ac{BNPL} providers apply is the mechanism under study. Imposing a genuine affordability test on
\ac{BNPL} is an intervention, and appears as the second lever of Submodel~15.

\textbf{A purely multiplicative peer-influence rule (Submodel~11).} Every household opens with no
\ac{BNPL} balance, so the reference-group adoption share is identically zero at initialisation and
adoption could never begin. The additive form of Equation~\ref{eq:peer} gives $q_{\text{base}}$ and
$\beta$ the roles that $p$ and $q$ play in Bass diffusion.

\subsection{Environment and Scheduling}

\textbf{Macroeconomic dynamics (Submodel~16).} Both Cardaci \cite{cardaci2018inequality} and
Madeira \cite{madeira2018chile} carry macro dynamics because both ask macro questions. Any
time-variation here would confound the effect of injecting \ac{BNPL} with a macro effect, making it
impossible to attribute a change in default to the injection.

\textbf{Uniform activation order (Submodel~17).} Credit in this model is scarce, rationed by the
affordability gate and by per-platform limits, so a fixed order would hand the same households
first claim on that scarce resource every tick. Comer and Loerch \cite{comer2013activation} find
significant differences in emergent behaviour across activation regimes, and Alizadeh and
Cioffi-Revilla \cite{alizadeh2015activation} find that no scheme dominates, so the choice is
reported and tested instead of assumed away.

